\subsection{\texorpdfstring{星代数 \(\mathrm{L}^{\infty}(\mathrm{X},\mu)\)}{星代数 \textbackslash mathrm\{L\}\^{}\{\textbackslash infty\}(\textbackslash mathrm\{X\},\textbackslash mu)}}

设 X 是局部紧拓扑空间，\(\mu\) 是 X 上的正测度。我们考虑交换星代数 \(\mathrm{L}^{\infty}(\mathrm{X},\mu)\) （I, p. 103 的例 4）。

引理 6. --- 设 \(g \in \mathcal{L}^{\infty}(\mathrm{X}, \mu)\)。\(g\) 在 \(\mathrm{L}^{\infty}(\mathrm{X}, \mu)\) 中的类的谱等于 \(\mu\)-本质像（\(g\) 的）。

用 \(\widetilde{g}\) 表示 \(g\) 在 \(\mathrm{L}^{\infty}(\mathrm{X},\mu)\) 中的类，用 \(S\) 表示 \(\mu\)-本质像（\(g\) 的）。设 \(z\in \mathbf{C} - \mathbf{S}\)。按定义，存在开邻域 \(U\) （\(z\) 的），使得 \(\mathrm{Y} = g^{-1}(\mathrm{U})\) 局部 \(\mu\)-零测。函数 \(h\) ——由 \(h(x) = (g(x) - z)^{-1}\) （当 \(x\notin \mathrm{Y}\)）与 \(h(x) = 0\) （当 \(x\in \mathrm{Y}\)）定义——属于 \(\mathcal{L}^{\infty}(\mathrm{X},\mu)\)。它的类 \(\widetilde{h}\) 在 \(\mathrm{L}^{\infty}(\mathrm{X},\mu)\) 中满足 \((\widetilde{g} -z)\widetilde{h} = 1\) ，因为 \((g(x) - z)h(x) = 1\) 对每个 \(x\) （只要它不属于局部 \(\mu\)-零测集 Y）成立。因此 \(z\in \mathbf{C} - \mathrm{Sp}(\widetilde{g})\) 。

反之，设 \(z \in \mathbf{C} - \mathrm{Sp}(\widetilde{g})\)。设 \(h \in \mathcal{L}^{\infty}(\mathrm{X},\mu)\) 是这样的函数：它的类是 \(\widetilde{g} - z\) 在 \(\mathrm{L}^{\infty}(\mathrm{X},\mu)\) 中的逆。存在实数 \(\mathrm{M} > 0\) ，使得 \(|h(x)| \leqslant \mathrm{M}\) 局部 \(\mu\)-几乎处处成立，此外我们有 \((g(x) - z)h(x) = 1\) 局部 \(\mu\)-几乎处处成立。设 U 是以 \(z\) 为心、\(\mathrm{M}^{-1}\) 为半径的开球（在 \(\mathbf{C}\) 中）；则 \(g^{-1}(\mathrm{U})\) 含于局部 \(\mu\)-零测集

\[
h ^ { - 1 } ( ] \mathrm{M} , + \infty [ ) \cup \{ x \in \mathrm{X} \mid ( g ( x ) - z ) h ( x ) \neq 1 \} ,
\]

从而 \(z \in \mathbf{C} - \mathrm{S}\) 。引理得证。

命题 4. --- 设 \(g \in \mathcal{L}^{\infty}(\mathrm{X},\mu)\) 。用 \(\widetilde{g}\) 表示 \(g\) 在 \(\mathrm{L}^{\infty}(\mathrm{X},\mu)\) 中的类，用 \(\mathrm{S}\) 表示 \(\widetilde{g}\) 的谱。映射 \(f \mapsto f \circ g\) （p. 184 的注）是从 \(\mathcal{C}(\mathrm{S})\) 到 \(\mathrm{L}^{\infty}(\mathrm{X},\mu)\) 中的与 \(\widetilde{g}\) 相伴的泛函演算映射。

设 \(f \in \mathcal{C}(\mathrm{S})\) 。函数 \(h = f \circ g\) 是 \(\mu\)-可测的（IV, p. 184 的引理 5）且有界。用 \(\widetilde{f \circ g}\) 表示它在 \(\mathrm{L}^{\infty}(\mathrm{X},\mu)\) 中的类。映射 \(f \mapsto \widetilde{f \circ g}\) 是从 \(\mathcal{C}(\mathrm{S})\) 到 \(\mathrm{L}^{\infty}(\mathrm{X},\mu)\) 中的连续有单位元的对合代数同态。由于 \(g(x)\) 属于 \(\mathrm{S}\) （局部 \(\mu\)-几乎处处；IV, p. 181 的引理 1 及引理 6），这个同态把 \(\mathrm{S}\) 上的恒等映射映到类 \(\widetilde{g}\) （函数 \(g\) 的类）。因此它与 \(\widetilde{g}\) 的泛函演算映射相合（I, p. 111 的命题 7）。
% semantic-fragments: legacy-input-seam
\relax{}
